\begin{pastexam}[label={px:ntut-113-ee-ds-08}]{北科大電機所（戊組）113 學年度 第 II.4 題}
（簡答；4 分）

For an ordered list of 11 items, how many comparisons will a linear search need before it knows
that a value is NOT in the list? (4\%)
\end{pastexam}

\begin{pastsoln}
答案是 $11$ 次（最壞情況）。

題目沒有給要找的值，也沒有說是哪一種線性搜尋，所以要分開講，取題目要的「最多要比幾次才確定
不在」。

\textbf{一般的線性搜尋}（不利用排序）：要確定某個值不在串列裡，必須和每一個元素都比過，
$11$ 個元素就是 $11$ 次。 不管要找的值是多少都一樣。

\textbf{利用排序提早停}：串列由小到大排好時，掃到第一個比目標大的元素就可以停，因為後面的
元素更大。 這時次數看目標落在哪裡：目標比第 $1$ 個元素小，比 $1$ 次就知道；目標落在第 $i - 1$
和第 $i$ 個元素之間，要比 $i$ 次；目標比最後一個還大，仍然要比完 $11$ 次。 所以最好 $1$ 次、
最壞 $11$ 次。 若目標等可能落在 $12$ 個空隙中的任何一個，平均是
\[
\frac{1 + 2 + \cdots + 11 + 11}{12} = \frac{66 + 11}{12} = \frac{77}{12} \approx 6.4
\]
次，大約 $n/2$（資結書第~1.8~節最好、最壞、平均情況的分法）。

題目問「需要幾次」而沒有指定目標，答最壞情況的 $11$ 次；有排序可利用時，可以補充平均大約
一半。
\end{pastsoln}
